\section{Overview}
Most of this part is dedicated to describing the features of various languages that are
relavent to this paper, in terms defined by this paper.
Each of these sections will start with a summary of the features provided by the language.
Next, it will explain the background knowledge necessary in order to understand the
feature descriptions.
This explanation will include an example of a "simple function and call" which serves two
purposes.
First, it provides the reader with a concrete example of what analysis looks like in the
context of this language.
For example, in Python "analysis" means assessing the behavior of the interpreter and the
bytocode emitted by the compiler.
Second, it provides a point of comparison when describing features.
In the Python analysis, the description of destructuring discusses how the bytecode for
this feature differs from a call that does not include destructuring.
This keeps each section shorter because they do not have to explain the baseline they are
being compared to.
It also makes it easier for a reader to focus only on the features that interest them
because they can read the simple call section then the feature section without having to
pick through other features to find the background information that they require.
Finally, each language analysis will describe each of the features.

Each feature description will start either with a statement that the feature is not
provided or a statement of how the feature is provided.
For example, in Python positional parameters are the "default behavior".
Tutorials (including the official tutorial contained in the repository) commonly introduce
functions by using positional parameters and introduce "keyword arguments" as a separate
feature at a later point.
Therefore, the Python section about positional values starts with "provided by default"
while the section about named values starts with "provided through keyword arguments".
This helps the reader understand the feature's relationship to the language and clarifies
what terminology they should expect to see if they are reading language documentation, or
what terms they should use if they want to perform a digital search for more information.
After this statement, any distinctive qualities of the feature as provided by the language
will be noted.

Finally, there will be 3 subsections: Syntax, Implementation, and Historical Record.
All 3 of these sections inform the implementation of the mechanism and perform a service
for the reader
\footnote{It should go without saying that the author is also a reader, albeit a
particularly invested one. =)}.

The syntax section explains what the source code looks like when the feature is used.
This helps readers who are unfamiliar with the language in question understand the code
snippets in the following subsections.

The implementation section explains the language behavior which causes the feature to be
provided.
This helps clarify the "concrete things" that the argument is associating.

The historical record section discusses, where possible, the motivation for the feature
and lessons learned from implementation and community response.
This gives the reader context about the environment the feature exists in, deepening their
understanding.

The final section synthesizes the information from the descriptions.
When different communities have similar concerns it will merge these concerns into a
single description.
When the concerns are distinct it will clarify the distinction.
It will also look for opportunities to apply solutions created by one community to a
concern raised by another.
Finally, it will provide a compact table listing all concerns.
The mechanism will either address each of these concerns or provide justifications for
leaving specific concerns unaddressed.
